% Part: incompleteness
% Chapter: representability-in-q
% Section: c-representable

\documentclass[../../../include/open-logic-section]{subfiles}

\begin{document}

\olfileid{inc}{req}{cre}
\olsection{$C$ وچلے فنکشن $\Th{Q}$ وچ قابلِ نمائندگی نیں}

سانوں وکھاؤنا اے کہ $C$ وچلا ہر فنکشن $\Th{Q}$ وچ قابلِ نمائندگی اے۔ آخر سانوں وکھاؤنا پئے گا کہ ہر $k$-دلیلی فنکشن $f(x_0,\dots,x_{k-1})$، جیہڑا $C$ وچ اے، نال فارمولا $!A_f(x_0,\dots,x_{k-1},y)$ کِویں جوڑیے جیہڑا اوہدی نمائندگی کرے۔

\begin{lem}
قدرتی عدد $n$ تے $m$ دِتّے ہون، تاں جے $n \neq m$، فیر $Q \Proves \num
n \neq \num m$۔
\end{lem}

\begin{proof}
$n$ اُتے استقرا نال وکھاؤ کہ ہر $m$ لئی، جے $n \neq m$، تاں
$Q \Proves \eq/[\num n][\num m]$۔

مُڈھلی صورت وچ $n = 0$۔ جے $m$، $0$ دے برابر نہیں، تاں $m = k +
1$ اے، کسے قدرتی عدد $k$ لئی۔ ساڈے کول مسلّمہ $\lforall[x][\eq/[0][x']]$ اے۔ مسوّر دے اک مسلّمے نال، $x$ دی تھاں $\num k$ رکھ کے، اسیں $\eq/[0][\num k']$ کڈھ سکدے آں۔ پر $\num k'$ تاں $\num m$ ای اے۔

استقرا دے قدم وچ اسیں منّ سکدے آں کہ دعویٰ $n$ لئی سچ اے، تے $n+1$ نوں ویکھدے آں۔ کوئی وی قدرتی عدد $m$ لو۔ دو صورتاں نیں۔ یا $m = 0$، یا کسے $k$ لئی $m = k+1$۔ پہلی صورت اُتلے طریقے نال ہو جاندی اے۔ دوجی صورت وچ منّ لو $n+1 \neq
k+1$۔ تاں $n \neq k$۔ $n$ لئی استقرا دے مفروضے نال
$\Th{Q} \Proves \eq/[\num n][\num k]$۔ ساڈے کول مسلّمہ اے
$\lforall[x][\lforall[y][\eq[x'][y'] \lif \eq[x][y]]]$۔ مسوّر دے اک مسلّمے نال $\eq[\num n'][\num k'] \lif \eq[\num n][\num
  k]$ ملدا اے۔ قضیائی منطق نال $\Th{Q}$ وچ
$\eq/[\num n][\num k] \lif \eq/[\num n'][\num k']$ کڈھ سکدے آں۔ اثباتِ مقدم نال $\eq/[\num n'][\num k']$ کڈھ سکدے آں۔ سانوں ایہو چاہیدا سی، کیوں جے $\num k'$ دراصل $\num m$ اے۔
\end{proof}

دھیان دیو، معاون قضیہ کوئی وڈی گل نہیں کردا۔ ایہ بنیادی طور اُتے کہندا اے کہ $\Th{Q}$ ثابت کر سکدا اے کہ وکھ وکھ عددی عبارتاں وکھ وکھ شےواں نوں دسدیاں نیں۔ مثلاً $\Th{Q}$، $0'' \neq 0'''$ ثابت کردا اے۔ پر ایہ عام طور اُتے وکھاؤن لئی کجھ دھیان چاہیدا اے۔ ایہ وی دھیان دیو کہ بھاویں اسیں استقرا ورتدے آں، ایہ استقرا \emph{باہر} اے، $\Th{Q}$ توں۔

اسیں صفر، جانشین، جمع، ضرب، برابری دے امتیازی فنکشن تے پروجیکشن فنکشاں دی نمائندگی کر سکاں گے۔ ہر صورت وچ نمائندگی کرن والا ڈھکدا فنکشن بالکل سدھا اے۔ مثلاً صفر دی نمائندگی ایہ فارمولا کردا اے:
\[
y = 0,
\]
جانشین دی نمائندگی ایہ فارمولا کردا اے:
\[
x_0' = y,
\]
تے جمع دی نمائندگی ایہ فارمولا کردا اے:
\[
x_0 + x_1 = y.
\]
اصل کم ایہ وکھاؤنا اے کہ $\Th{Q}$ متعلقہ بیان ثابت کر سکدا اے۔ مثلاً ایہ آکھن لئی کہ اُتلا فارمولا جمع دی نمائندگی کردا اے، سانوں وکھاؤنا پئے گا کہ قدرتی عدداں $m$ تے $n$ دے ہر جوڑے لئی $\Th{Q}$ ایہ ثابت کردا اے:
\[
\eq[\num n + \num m][\num {n+m}]
\]
تے
\[
\lforall[y][(\eq[(\num n + \num m)][y] \lif \eq[y][\num {n+m}])].
\]

ترکیب بارے کیہ؟ منّ لو $h$ دی تعریف ایویں اے:
\[
h(x_0,\dots,x_{l-1}) = f(g_0(x_0,\dots,x_{l-1}), \dots,
g_{k-1}(x_0,\dots,x_{l-1})).
\]
ایتھے اسیں پہلاں ای فارمولے $!A_f, !A_{g_0}, \dots,
!A_{g_{k-1}}$ لبھ چکے آں جیہڑے بالترتیب فنکشاں $f,g_0,\dots,g_{k-1}$ دی نمائندگی کردے نیں۔ فیر اسیں فارمولا $!A_h$، جیہڑا $h$ دی نمائندگی کردا اے، ایویں تعریف کر سکدے آں کہ $!A_h(x_0,\dots,x_{l-1},y)$ ایہ ہووے:
\begin{multline*}
  \lexists[z_0\dots][\lexists[z_{k-1}][(!A_{g_0}(x_0,\dots,x_{l-1},z_0) \land
  \dots \land !A_{g_{k-1}}(x_0,\dots,x_{l-1},z_{k-1}) \land]]\\
  !A_f(z_0,\dots,z_{k-1},y)).
\end{multline*}

آخر غیر محدود تلاش ویکھو۔ منّ لو $g(x,\vec z)$ باقاعدہ اے تے $\Th{Q}$ وچ قابلِ نمائندگی اے، مثلاً فارمولا $!A_g(x,\vec
z,y)$ اوہدی نمائندگی کردا اے۔ $f$ دی تعریف $f(\vec z) = \mu x \; g(x,\vec z)$ نال کرو۔ اسیں فارمولا $!A_f(\vec z,y)$ لبھنا چاہندے آں جیہڑا $f$ دی نمائندگی کرے۔ اک قدرتی چُون ایہ اے:
\[
!A_f(\vec z,y) \ident !A_g(y,\vec z,0) \land \lforall[w][(w < y
\lif \lnot !A_g(w,\vec z,0))].
\]

کجھ ہور معاون قضیاں نال وکھایا جا سکدا اے کہ ایہ کم کردا اے، مثلاً ایہناں نال۔

\begin{lem}
  ہر متغیر $x$ تے ہر قدرتی عدد $n$ لئی، $\Th{Q}$،
  $\eq[(x'
  + \num n)][(x + \num n)']$ ثابت کردا اے۔
\end{lem}

فیر دَس دیئے کہ ایہ ایس گل توں کمزور اے کہ $\Th{Q}$،
$\lforall[x][\lforall[y][(x' + y) = (x +y)']]$ ثابت کردا اے (جیہڑی غلط اے)۔

\begin{proof}
ثبوت، جیویں عام ورتدے آں، $n$ اُتے استقرا نال اے۔ مُڈھلی صورت وچ $n =
0$، سانوں وکھاؤنا اے کہ $\Th{Q}$، $\eq[(x'+0)][(x+0)']$ ثابت کردا اے۔ پر ساڈے کول ایہ اے:
\begin{eqnarray*}
& & \eq[(x' + 0)][x'] \quad \text{مسلّمہ 4 توں} \\
& & \eq[(x + 0)][x] \quad \text{مسلّمہ 4 توں} \\
& & \eq[(x+0)'][x'] \quad \text{سطر 2 نال} \\
& & \eq[(x' + 0)][(x + 0)'] \quad \text{سطراں 1 تے 3}
\end{eqnarray*}
استقرا دے قدم وچ منّ سکدے آں کہ اسیں $\eq[(x' +
  \num n)][(x + \num n)']$، $\Th{Q}$ وچ اخذ کر چکے آں۔ چونکہ $\num{n+1}$،
$\num
n'$ اے، سانوں وکھاؤنا اے کہ $\Th{Q}$، $\eq[(x' + \num n')][(x + \num
n')']$ ثابت کردا اے۔ ایہ برابریاں دی لڑی $\Th{Q}$ وچ اخذ ہو سکدی اے:
\begin{eqnarray*}
(x' + \num n') & \eq & (x' + \num n)' \quad \text{مسلّمہ 5} \\
& \eq & (x + \num n')' \quad \text{استقرا دے مفروضے توں}
\end{eqnarray*}
\end{proof}

\begin{lem}
\begin{enumerate}
\item $\Th{Q}$، $\lnot (x < \num 0)$ ثابت کردا اے۔
\item ہر قدرتی عدد $n$ لئی، $\Th{Q}$ ایہ ثابت کردا اے:
\[
x < \num {n+1} \lif (\eq[x][\num 0] \lor \dots \lor \eq[x][\num n]).
\]
\end{enumerate}
\end{lem}

\begin{proof}
آؤ 1 تے 2 دا کجھ حصہ غیر رسمی طریقے نال کریئے (یعنی بس ایہ اشارے دیئے کہ رسمی !!{derivation} کِویں بنایا جا سکدا اے)۔

پہلے حصے لئی، $<$ دی تعریف مطابق، سانوں $\lnot
\lexists[y][\eq[(y'+x)][0]]$، $\Th{Q}$ وچ ثابت کرنا اے۔ پہلے درجے دی منطق دے مسلّمے تے قاعدے ورتیاں ایہ $\lforall[y][\eq/[(y' +
    x)][0]]$ دے برابر اے۔ خیال ایہ اے۔ منّ لو $\eq[(y'+x)][0]$۔ جے $x$ صفر اے، تاں $\eq[(y' + 0)][0]$۔ پر $\Th{Q}$ دے مسلّمہ 4 نال
$\eq[(y' + 0)][y']$ تے مسلّمہ 2 نال $\eq/[y'][0]$، سو تضاد اے۔ ایس لئی
$\lforall[y][\eq/[(y' +x)][0]]$۔ جے $x$، $0$ نہیں، تاں مسلّمہ 3 نال کوئی $z$ اے جیہدے لئی $\eq[x][z']$۔ پر فیر
$\eq[(y' + z')][0]$۔ مسلّمہ~5 نال $\eq[(y' + z)'][0]$، سو فیر مسلّمہ~2 نال تضاد اے۔

دوجے حصے لئی $n$ اُتے استقرا ورتو۔ مُڈھلی صورت ویکھو، جدوں $n =0$۔ اوس صورت وچ سانوں $x < \num 1 \lif \eq[x][\num
  0]$ وکھاؤنا اے۔ منّ لو $x < \num 1$۔ تاں $<$ دی تعریف والے مسلّمے نال $\lexists[y][\eq[(y'+x)][0']]$۔ منّ لو $y$ وچ ایہ خاصیت اے، سو $\eq[y'+x][0']$۔

سانوں $\eq[x][0]$ وکھاؤنا اے۔ مسلّمہ 3 نال جے $x$ صفر نہیں تاں $z'$ دے برابر اے، کسے $z$ لئی۔ فیر $\eq[(y' + z')][0']$۔ مسلّمہ~5، جیہڑا
$\Th{Q}$ دا اے، نال $\eq[(y'+z)'][0']$۔ مسلّمہ 1 نال $\eq[(y' +
    z)][0]$۔ پر تعریف مطابق ایہدا مطلب $z < 0$ اے، جیہڑا حصّہ~1 نال ٹکراؤندا اے۔
\end{proof}

\begin{explain}
اسیں وکھایا اے کہ قابلِ حساب فنکشاں دا سیٹ $\Th{Q}$ وچ قابلِ نمائندگی فنکشاں دا سیٹ اے۔ اصل وچ ثبوت ہور عام اے۔ نمائندگی دی تعریف توں سوکھا دِسدا اے کہ ہر نظریہ جیہڑا $\Th{Q}$ دا وادھا اے (یا جیہدے وچ $\Th{Q}$ دی تعبیر کیتی جا سکدی اے)، قابلِ حساب فنکشاں دی نمائندگی کر سکدا اے۔ اُلٹے پاسے، ہر !!{derivation} نظام وچ جیہدے وچ ثبوت دا تصور قابلِ حساب اے، ہر قابلِ نمائندگی فنکشن قابلِ حساب اے۔ سو، مثلاً، قابلِ حساب فنکشاں دا سیٹ پیانو حساب، یا زرمَیلو–فرینکل سیٹ نظریے، وچ قابلِ نمائندگی فنکشاں دا سیٹ وی اے۔ جیویں گوڈل نے دَسیا، ایہ کجھ حیرانی والی گل اے۔ اسیں ویکھاں گے کہ ثابت پذیری دے سوال ایس گل اُتے بہت نربھر کردے نیں کہ تسی کیہڑا نظریہ ویکھ رہے او۔ موٹے طور اُتے، مسلّمے جِنّے تکڑے ہون، اُنا ودھ ثابت کر سکدے او۔ پر مسلّماتی نظرییاں دی اک وڈی حد وچ قابلِ نمائندگی فنکشن عین قابلِ حساب فنکشن ای نیں۔
\end{explain}

\end{document}
